%% FILE 07 | CHỦ SỞ HỮU: Người 3 | HẠN: 13/10
%% File này MỞ section 7. Người 2 viết tiếp phần quyết định ở file 08 (subsection).
%% EXPECTED: tóm tắt từng finding gắn với assumption nào (A1, A2...), mức độ tin cậy, giới hạn.
%% Label tiền tố p3:
\section{Kết quả khảo sát và quyết định sản phẩm}
\subsection{Những gì tìm thấy}
Khảo sát tự báo cáo được thực hiện từ 03/10 đến 12/10/2026; trong 38 dòng dữ liệu thô, sau khi loại email trùng và giữ phản hồi gửi muộn nhất còn n=35. Có 21/35 người từng lên phòng, 14/35 chưa từng lên phòng và 17/35 không lên phòng lần nào trong tháng gần đây. Đây là mẫu thuận tiện, không đại diện cho toàn bộ thành viên CLB; thông tin về kênh phát (nhóm nội bộ và group sinh viên) được nêu tại mục 5.1; do thu thập qua đường dẫn mở nên nhóm không kiểm soát tỷ lệ phân bổ giữa người thường lên phòng và người chưa từng lên phòng. Các câu trả lời là kinh nghiệm hoặc ý định tự báo cáo, chưa phải hành vi quan sát; câu hỏi cho phép chọn nhiều đáp án nên tổng lựa chọn có thể vượt 35.
\par
\setlength{\tabcolsep}{4pt}
\setlength{\extrarowheight}{1pt}
{\centering
\begin{tabularx}{0.94\textwidth}{|>{\raggedright\arraybackslash}p{0.07\textwidth}|>{\raggedright\arraybackslash}p{0.45\textwidth}|X|}
\hline
\textbf{Giả định} & \textbf{Kết quả khảo sát} & \textbf{Độ tin cậy / giới hạn} \\ \hline
A1 & Trong khảo sát tự báo cáo, mẫu thuận tiện (n=35), 21/35 xem Google Sheet, Form hoặc lịch trực để biết ai đang trực/phòng có ai; 6/35 không kiểm tra trước, 3/35 ghé phòng và 3/35 chưa có cách nào tiện. Trong nhóm từng lên phòng (n=21), 14/21 xem Sheet/Form/lịch. Về trải nghiệm, 16/35 chưa gặp vấn đề đáng kể (10/21 trong nhóm từng lên phòng), 12/35 từng gặp lịch thay đổi mà mọi người không biết (7/21), và 8/35 không biết ai đang giữ hoặc trực phòng (5/21). Tín hiệu nhu cầu không đồng nhất và chưa xác nhận mức độ phổ biến trong toàn bộ thành viên. & Độ tin cậy trung bình đối với nhu cầu tự báo cáo. Khảo sát không quan sát hành vi thực tế; một số người chưa từng sử dụng phòng và có thể chưa nắm rõ quy trình hiện tại. \\ \hline
A2 & Trong khảo sát tự báo cáo, mẫu thuận tiện (n=35), 24/35 nói sẽ check-in/out mỗi lần nếu thao tác nhanh; trong nhóm từng lên phòng, con số là 17/21. Ngoài ra, 5/35 chỉ khi có ca trực, 3/35 khi thấy lợi ích rõ ràng và 3/35 chưa chắc, cần dùng thử. Đây là mức sẵn sàng tự báo cáo có điều kiện, chưa chứng minh hành vi sử dụng lặp lại. & Độ tin cậy thấp đến trung bình. Câu trả lời thể hiện ý định trong tình huống giả định, chưa phải hành vi quan sát được khi sử dụng sản phẩm. \\ \hline
A3 & Trong khảo sát tự báo cáo, mẫu thuận tiện (n=35), phương án được chọn nhiều nhất là máy đặt tại phòng quét thẻ SV (13/35); trong nhóm từng lên phòng, 9/21 chọn cách này. Có 10/35 chọn mở app cá nhân quét QR của phòng (5/21 trong nhóm từng lên phòng), 5/35 chọn mở app bấm nút và khoảng 5/35 chọn cách tự động hơn. Khảo sát không hỏi trực tiếp mô hình người đang trực dùng điện thoại quét thẻ/QR của thành viên. Vì vậy mô hình này chỉ là phương án thử nghiệm, chưa có bằng chứng ủng hộ trực tiếp. & Độ tin cậy trung bình về sự đa dạng sở thích; phương án máy cố định được chọn nhiều nhất trong các lựa chọn được nêu, nhưng khảo sát không kiểm tra khả năng sử dụng thành công của các phương thức. \\ \hline
\end{tabularx}\par}
\clearpage
{\centering
\begin{tabularx}{0.94\textwidth}{|>{\raggedright\arraybackslash}p{0.07\textwidth}|>{\raggedright\arraybackslash}p{0.45\textwidth}|X|}
\hline
\textbf{Giả định} & \textbf{Kết quả khảo sát} & \textbf{Độ tin cậy / giới hạn} \\ \hline
A4 & Trong khảo sát tự báo cáo, lấy mẫu thuận tiện (n=35), 21/35 cho biết có thể tham gia trực nếu lịch phù hợp, 10/35 chưa chắc, 3/35 không có nhu cầu và 1/35 trả lời có. Kết quả gợi ý lịch linh hoạt có thể ảnh hưởng đến ý định tham gia, nhưng chưa xác nhận tỷ lệ đăng ký hoặc hoàn thành ca thực tế. & Độ tin cậy thấp đến trung bình. Đây là ý định tự báo cáo; khảo sát chưa kiểm chứng việc đăng ký và hoàn thành ca trực. \\ \hline
A5 & Trong khảo sát tự báo cáo, lấy mẫu thuận tiện (n=35), ở câu hỏi chọn vài việc nên làm tốt nhất, 21/35 chọn biết phòng đang mở hay đóng, 16/35 chọn biết hiện tại phòng có ai và 16/35 chọn đăng ký/đổi ca trực; 9/35 chọn xem lịch trực và 9/35 chọn xem ai đang trực. Ngoài ra, 12/35 từng gặp lịch thay đổi mà mọi người không biết. Các lựa chọn cho thấy có quan tâm đến đăng ký/đổi ca, song không chứng minh hành vi hoặc quy trình khả thi. & Độ tin cậy trung bình đối với các nhu cầu tự báo cáo; chưa thể suy ra hành vi sử dụng hoặc quy trình phù hợp của CLB. \\ \hline
A6 & Trong khảo sát tự báo cáo, lấy mẫu thuận tiện (n=35), 13/35 nêu dễ quên, 8/35 mạng trường không ổn định, 7/35 vội, 6/35 phải mở app/link riêng, 6/35 quá nhiều bước và 1/35 không thấy lợi ích trực tiếp là điều dễ khiến họ bỏ qua check-in/out; 16/35 cho biết thao tác là bình thường, không phiền. Điều này hỗ trợ thiết kế ít thao tác và báo lỗi/thử lại khi mất mạng, nhưng chưa chứng minh nhắc nhở hoặc offline sẽ giải quyết rào cản. & Độ tin cậy trung bình đối với các rào cản tự báo cáo; hiệu quả của giải pháp cần được kiểm chứng bằng thử nghiệm sử dụng thực tế. \\ \hline
A7 & Khảo sát tập trung vào thành viên và không hỏi trực tiếp ban quản lý về thời gian, quy trình hoặc khó khăn khi tổng hợp dữ liệu. Do đó, chưa có bằng chứng trực tiếp để xác nhận nhu cầu giảm công sức quản trị hoặc nhu cầu về báo cáo tập trung. & Chưa đủ bằng chứng để kết luận. Cần phỏng vấn hoặc quan sát ban quản lý CLB và đối chiếu quy trình tổng hợp hiện tại trước khi xác định ưu tiên cho dashboard hoặc báo cáo. \\ \hline
\end{tabularx}\par}
\setlength{\tabcolsep}{6pt}
\setlength{\extrarowheight}{0pt}
